% Generated draft: requires mathematical and source review.
\documentclass[11pt]{article}
\usepackage[margin=25mm]{geometry}
\usepackage{fontspec}
\setmainfont{Latin Modern Roman}
\usepackage{amsmath,amssymb,mathtools,unicode-math}
\setmathfont{Latin Modern Math}
\usepackage{enumitem,needspace,xurl,hyperref}
\setlength{\parindent}{0pt}
\setlength{\parskip}{5pt}
\setlength{\emergencystretch}{3em}
\newcommand{\N}{\mathbb{N}}
\newcommand{\Z}{\mathbb{Z}}
\newcommand{\Q}{\mathbb{Q}}
\newcommand{\R}{\mathbb{R}}
\newcommand{\C}{\mathbb{C}}
\newcommand{\F}{\mathbb{F}}
\newcommand{\A}{\mathbb{A}}
\newcommand{\PP}{\mathbb{P}}
\newcommand{\E}{\mathbb{E}}
\newcommand{\eps}{\varepsilon}
\newcommand{\dd}{\,\mathrm{d}}
\begin{document}
\section*{UnsolvedMath ID: 700769}
\textbf{Source problem code:} catalogue entry
\subsection{Definitions and mathematical statement}
Let \(\Pi\) be a finite projective plane of order \(n\geq 2\), and let \(\ell_\infty\) be a distinguished line. Thus each line has \(n+1\) points, any two distinct points lie on a unique line, any two distinct lines meet in a unique point, and the plane is nondegenerate. The points outside \(\ell_\infty\) are called affine points.

A collineation is an incidence-preserving permutation of the points and lines. A translation with respect to \(\ell_\infty\) is either the identity or an elation with axis \(\ell_\infty\): a nonidentity collineation fixing \(\ell_\infty\) pointwise and fixing every line through some point \(C\in\ell_\infty\) setwise. The pair \((\Pi,\ell_\infty)\) is a translation plane if its translation group acts regularly on the affine points; that is, for any affine points \(P,Q\), exactly one translation sends \(P\) to \(Q\).

Write
\[
G=\{g\in\operatorname{Aut}(\Pi):g(\ell_\infty)=\ell_\infty\}.
\]
The action of \(G\) on the points of \(\ell_\infty\) is doubly transitive if, for every \(A,B,A',B'\in\ell_\infty\) with \(A\ne B\) and \(A'\ne B'\), there exists \(g\in G\) such that
\[
g(A)=A',\qquad g(B)=B'.
\]

\textbf{Classification problem (catalogue entry).} Describe all finite translation planes \((\Pi,\ell_\infty)\) for which this action is doubly transitive.

Here the collineation group is interpreted as the full collineation group of the affine plane, equivalently the setwise stabilizer of its line at infinity in the projective completion. Existence of a subgroup acting doubly transitively on that line is equivalent to double transitivity of this full stabilizer. The source does not explicitly specify the equivalence relation for the requested classification; classification up to incidence isomorphisms carrying the distinguished line to the distinguished line is the natural marked-plane formulation used here.
\subsection{Short English statement}
Describe the finite translation planes whose collineations act transitively on ordered pairs of distinct points at infinity. N. D. Podufalov proposed this as a ``well-known problem.'' The supplied registry labels it open; that label has not been independently verified.
\subsection{Sources}
\begin{itemize}
\item Supplied local UnsolvedMath registry record: permanent ID \(700769\), exact source problem code \texttt{catalogue entry}, title \emph{Kourovka Notebook Problem 11.77}. The supplied status is \texttt{open}.
\item \emph{The Kourovka Notebook (Unsolved Problems in Group Theory)}, Problem 11.77, originating in the 11th issue (1990); the supplied background identifies the 21st issue edition (2026), arXiv identifier \texttt{1401.0300v48}. The registry describes the Issues 1--21 collection as edited by E. I. Khukhro and V. D. Mazurov.
\item Proposed by N. D. Podufalov, 1990. The supplied statement includes the qualification ``Well-known problem.''
\item The references and attribution above are supplied source data and have not been independently checked. No separate literature assessment or literature-verification date is supplied.
\end{itemize}
\end{document}
